\providecommand{\rank}{\operatorname{r}}
\providecommand{\RR}{\mathbb R}
\providecommand{\CC}{\mathbb C}
\providecommand{\KK}{\mathbb K}

% ---------- Exercise 78 ----------
\setcounter{enumi}{77}
\item {\bfseries\boldmath
For each of the following systems of vectors, determine whether it is linearly independent in the indicated vector space. If it is dependent, express one vector as a linear combination of the others. Using these conclusions, determine a basis of the vector subspace generated by the given system.
}

\begin{enumerate}[label=\alph*)]

\item $\bigl((1,2),(-2,-4)\bigr)$ in $\RR^2$.

\begin{answer}
A linearly independent generating system is a basis of its span.
In each part the vectors are kept in the order in which they are listed.

For $v_1=(1,2)$, $v_2=(-2,-4)$, we have $v_2=-2v_1$.
The system is dependent, and $\bigl((1,2)\bigr)$ is a basis of its span.
\end{answer}

\item $\bigl((1,1),(-1,-2)\bigr)$ in $\RR^2$.

\begin{answer}
The system $\bigl((1,1),(-1,-2)\bigr)$ is independent, since
$\det\begin{bmatrix}1&-1\\1&-2\end{bmatrix}=-1\ne0$.
It is itself a basis of its span, which is $\RR^2$.
\end{answer}

\item $\bigl((1,2)\bigr)$ in $\RR^2$.

\begin{answer}
The single vector $(1,2)$ is nonzero, so the system is independent
and is a basis of its one-dimensional span.
\end{answer}

\item $\bigl((2,-1,2),(2,2,2),(4,1,4)\bigr)$ in $\RR^3$.

\begin{answer}
Put $v_1=(2,-1,2)$, $v_2=(2,2,2)$, $v_3=(4,1,4)$.
Then $v_3=v_1+v_2$, so the system is dependent.
The first two vectors are independent: $\alpha v_1+\beta v_2=0$ gives
$\alpha+\beta=0$ and $-\alpha+2\beta=0$, hence $\alpha=\beta=0$.
A basis is $\bigl((2,-1,2),(2,2,2)\bigr)$.
\end{answer}

\item $\bigl((2,-1,2),(2,2,2),(1,1,-1)\bigr)$ in $\RR^3$.

\begin{answer}
For $\bigl((2,-1,2),(2,2,2),(1,1,-1)\bigr)$, the determinant of the
matrix of columns is
\[
 \det\begin{bmatrix}2&2&1\\-1&2&1\\2&2&-1\end{bmatrix}=-12\ne0.
\]
The system is independent and is itself a basis of $\RR^3$.
\end{answer}

\item $\bigl((0,1,2,0),(2,0,1,1),(-2,4,8,-1)\bigr)$ in $\RR^4$.

\begin{answer}
For $v_1=(0,1,2,0)$, $v_2=(2,0,1,1)$, $v_3=(-2,4,8,-1)$,
suppose $\alpha v_1+\beta v_2+\gamma v_3=0$.
Coordinates $1,2$ give $\beta=\gamma$, $\alpha=-4\gamma$.
Coordinate $3$ then gives $2\alpha+\beta+8\gamma=\gamma=0$.
All coefficients vanish, so the three given vectors are independent
and form a basis of their span.
\end{answer}

\item $\left(\begin{bmatrix}2&-1\\1&1\end{bmatrix},\begin{bmatrix}0&1\\1&0\end{bmatrix},\begin{bmatrix}2&2\\2&2\end{bmatrix}\right)$ in $M_2(\RR)$.

\begin{answer}
Let
\[
 A_1=\begin{bmatrix}2&-1\\1&1\end{bmatrix},\quad
 A_2=\begin{bmatrix}0&1\\1&0\end{bmatrix},\quad
 A_3=\begin{bmatrix}2&2\\2&2\end{bmatrix}.
\]
If $\alpha A_1+\beta A_2+\gamma A_3=0$, entries $(1,1)$ and $(2,2)$ give
$2\alpha+2\gamma=0$ and $\alpha+2\gamma=0$.
Thus $\alpha=\gamma=0$, and entry $(1,2)$ gives $\beta=0$.
The system is independent; $(A_1,A_2,A_3)$ is a basis of its span.
\end{answer}

\item $\left(\begin{bmatrix}4&2&2\\-2&0&2\end{bmatrix},\begin{bmatrix}-2&-1&1\\1&0&-1\end{bmatrix}\right)$ in $M_{2\times3}(\RR)$.

\begin{answer}
Let
\[
 A_1=\begin{bmatrix}4&2&2\\-2&0&2\end{bmatrix},\qquad
 A_2=\begin{bmatrix}-2&-1&1\\1&0&-1\end{bmatrix}.
\]
In $\alpha A_1+\beta A_2=0$, entries $(1,1)$ and $(1,3)$ give
$4\alpha-2\beta=0$, $2\alpha+\beta=0$.
Hence $\alpha=\beta=0$. The given pair is a basis of its span.
\end{answer}

\item $\bigl(2x,x^2+3x,x^2+3x+2\bigr)$ in $\RR_2[x]$.

\begin{answer}
For $p_1=2x$, $p_2=x^2+3x$, $p_3=x^2+3x+2$, compare coefficients
in $\alpha p_1+\beta p_2+\gamma p_3=0$.
The constant term gives $\gamma=0$, the $x^2$ coefficient gives $\beta=0$,
and the $x$ coefficient gives $\alpha=0$.
Thus $(p_1,p_2,p_3)$ is a basis of its span, namely $\RR_2[x]$.
\end{answer}

\item $\bigl(x^3-2x+1,\ 5,\ x^3+3x^2-1,\ -2x^3+3x^2-3x\bigr)$ in $\RR_3[x]$.

\begin{answer}
Put
\[
\begin{gathered}
 p_1=x^3-2x+1,\qquad p_2=5,\\
 p_3=x^3+3x^2-1,\qquad p_4=-2x^3+3x^2-3x.
\end{gathered}
\]
In $\alpha p_1+\beta p_2+\gamma p_3+\delta p_4=0$, the $x^2$ coefficient
gives $\gamma=-\delta$, and the $x^3$ coefficient then gives $\alpha=3\delta$.
The $x$ coefficient is $-2\alpha-3\delta=-9\delta=0$.
Hence $\delta=\alpha=\gamma=0$, and the constant term gives $5\beta=0$.
The four polynomials are independent and form a basis of $\RR_3[x]$.
\end{answer}
\end{enumerate}
